\section{Overview}
Living in a shared apartment or house involves numerous responsibilities that are shared among housemates, including tracking expenses, paying bills, buying groceries, and performing chores. \textit{HouseMate} is a software solution designed to streamline the management of these tasks. The application aims to reduce disputes and organizational overhead by providing a single, reliable point of truth shared between roommates.

\section{Statement of Problems and Solutions}

In many shared living environments, relying on group chats and personal notes to manage household commitments is the standard. While fast and immediate, this decentralized approach inevitably leads to unnecessary stress, wasted time, and avoidable conflicts among roommates. 

\textit{HouseMate} is designed to eliminate this friction. By digitalizing the household management ecosystem, it provides an automated, centralized platform that directly targets and resolves the most common pain points of sharing an apartment:

\begin{description}
    \item[1. Financial friction] \hfill \\
    \textit{The Problem:} Lost receipts, forgotten verbal agreements and the unequal division of bills create frustrating financial inconsistencies and tensions. \\
    \textit{The Solution:} \textbf{Centralized expense and debt tracking.} HouseMate acts as an impartial financial ledger. Residents insert transactions and the system logs all shared expenses and automatically calculates each roommate's individual debts and monthly contributions, ensuring absolute fairness without awkward conversations. 

    \item[2. Division of chores] \hfill \\
    \textit{The Problem:} The lack of a clear and visible schedule often results in an unequal distribution of household chores, leaving a few members to do most of the work. \\
    \textit{The Solution:} \textbf{Smart chore management.} This module allows flatmates to seamlessly define and assign specific tasks to each other. It tracks completion statuses in real time and automatically flags overdue chores, keeping the whole house accountable and the environment clean.

    \item[3. Grocery shopping] \hfill \\
    \textit{The Problem:} Disorganized communication leads to highly inefficient grocery runs where essential items are either forgotten entirely or accidentally purchased multiple times by multiple people. \\
    \textit{The Solution:} \textbf{Real-time shopping lists.} HouseMate provides dynamic, shared shopping lists that users can customize at will. Anyone at the store can see exactly what the house needs at all times, without guessing and wasting money.
\end{description}

By directly pairing these everyday hurdles with intuitive and automated solutions, HouseMate transforms a stressful shared apartment into a seamlessly managed home.

\section{Use of Artificial Intelligence}\label{sec:ai}

During the development of \textit{HouseMate}, Artificial Intelligence tools were utilized to assist with repetitive coding tasks and to generate the final styled versions of architectural diagrams based on our detailed structure previously built before and during development. It is important to emphasize that all AI-generated outputs were strictly supervised, reviewed, and refined under careful human control to ensure accuracy and strict alignment with the project requirements.